% !TeX root = ../main.tex
\chapter{Estado del arte}\label{cap:estado-arte}

Revisa críticamente lo que ya existe: no basta con enumerar trabajos,
hay que compararlos y detectar el hueco que cubre este proyecto. Frase añadida para probar la CI.

\section{Cómo citar}

Las referencias se guardan en \texttt{bibliografia/referencias.bib} y se
citan con \verb|\cite{clave}|: por ejemplo, la tipografía de este documento
se basa en \TeX{}~\cite{Knuth1984TeXbook} y \LaTeX{}~\cite{Lamport1994LaTeX}.
Se pueden citar varias a la vez~\cite{Knuth1984TeXbook, Lamport1994LaTeX}.
Usa claves del tipo \texttt{AutorAñoTema} e incluye siempre el DOI cuando
exista.

\section{Comparativa}

Una tabla comparativa suele ser la forma más clara de cerrar el capítulo
(\cref{tab:comparativa}).

\begin{table}[H]
  \centering
  \caption{Comparativa de soluciones existentes.}
  \label{tab:comparativa}
  \begin{tabular}{lccc}
    \toprule
    Solución & Coste & Precisión & Disponibilidad \\
    \midrule
    Alternativa A & Alto  & Alta  & Comercial \\
    Alternativa B & Bajo  & Media & Abierta   \\
    \textbf{Este trabajo} & \textbf{Bajo} & \textbf{Alta} & \textbf{Abierta} \\
    \bottomrule
  \end{tabular}
\end{table}

\section{Posicionamiento del trabajo}

Explica qué aporta este trabajo respecto a lo revisado.
